%% ma: L12-B05-0001
%% lop: 12
%% chuong: 1
%% bai: 5
%% dang: 12.5.1
%% loai: TLN
%% muc: VD
%% dap_an: "333"
%% nguon: "0. ĐỀ THAM KHẢO TN THPT 2025.pdf (D00), Phần III Câu 5"
%% kien_thuc: "Bài 5 (bài toán tối ưu thực tiễn bằng đạo hàm), Bài 2 (GTLN trên đoạn)"
%% trang_thai: cho-duyet
%% ly_do: "Dạng toán mới đề xuất 12.5.1, chờ giáo viên duyệt"
\begin{ex}
Một doanh nghiệp dự định sản xuất không quá $500$ sản phẩm. Nếu doanh nghiệp sản xuất $x$ sản phẩm ($1\le x\le500$) thì doanh thu nhận được khi bán hết số sản phẩm đó là $F(x)=x^3-1999x^2+1001000x+250000$ (đồng), trong khi chi phí sản xuất bình quân cho một sản phẩm là $G(x)=x+1000+\dfrac{250000}{x}$ (đồng). Doanh nghiệp cần sản xuất bao nhiêu sản phẩm để lợi nhuận thu được là lớn nhất?
\shortans{333}
\loigiai{Lợi nhuận $f(x)=F(x)-x\,G(x)=x^3-2000x^2+1000000x$, $x\in[1;500]$.
$f'(x)=3x^2-4000x+1000000=0\Leftrightarrow x=\frac{1000}{3}\approx333{,}3$ (loại $x=1000$). $f$ đồng biến trên $\left[1;\frac{1000}{3}\right]$, nghịch biến trên $\left[\frac{1000}{3};500\right]$.
Vì $x$ nguyên: $f(333)=148\,148\,037>f(334)=148\,147\,704$. Vậy cần sản xuất $333$ sản phẩm.}
\end{ex}
